\documentclass[11pt]{article}
\usepackage{assets/tex/notes/eric-notes-style}

\begin{document}

\notetitle{MATH 325: Complex Analysis Notes}{Midterm and Final Review}

\topic{Complex Numbers and Functions}

\entry{Complex number}
\[
\operatorname{Re} z=\frac{z+\overline z}{2},\qquad
\operatorname{Im} z=\frac{z-\overline z}{2i},\qquad
|z|^2=z\overline z.
\]
\[
\text{Euler's formula: } e^{i\theta}=\cos\theta+i\sin\theta \quad (\theta\in\mathbb R).
\]
\[
i^n=\begin{cases}
1, & n\equiv0\pmod4,\\
i, & n\equiv1\pmod4,\\
-1, & n\equiv2\pmod4,\\
-i, & n\equiv3\pmod4.
\end{cases}
\]

\entry{De Moivre's formula}
\[
\forall n\in\mathbb N,\qquad
(\cos\theta+i\sin\theta)^n=e^{in\theta}=\cos(n\theta)+i\sin(n\theta).
\]

\entry{Roots of unity}
\[
\text{For the } n\text{th roots of unity }z^n=1,\ z_k\in\mathbb C,
\]
\[
z_k=\exp\left(\frac{2\pi i k}{n}\right),\qquad k=0,1,\ldots,n-1.
\]

\entry{Complex exponential}
\[
\exp:\mathbb C\to\mathbb C,\qquad e^z=e^x(\cos y+i\sin y).
\]

\entry{Complex logarithm}
\[
\log z=\ln |z|+i\arg z.
\]

\entry{Complex powers}
\[
z^\alpha=e^{\alpha\log z}=e^{\alpha(\ln|z|+i\arg z)}.
\]
\[
a\in\mathbb R,\ a\ne0,\ n\in\mathbb Z,\quad \text{then }z^a
\text{ takes a distinct values.}
\]
\[
a\in\mathbb C \text{ and } a\notin\mathbb Z,\quad \text{then }z^a
\text{ takes infinitely many distinct values.}
\]

\entry{Trigonometric and hyperbolic functions}
\[
\cos z=\frac{e^{iz}+e^{-iz}}{2},\qquad
\sin z=\frac{e^{iz}-e^{-iz}}{2i}.
\]
\[
\cosh z=\frac{e^z+e^{-z}}{2},\qquad
\sinh z=\frac{e^z-e^{-z}}{2}.
\]

\entry{Images of exp}
\[
\text{Let } x=x_0 \text{ be constant: } e^z=e^x(\cos y+i\sin y),
\quad y\in\mathbb R.
\]
\[
\text{The image is a circle w/ radius } e^{x_0}.
\]
\[
\text{Let } y=y_0 \text{ be constant: } e^z=e^x(\cos y_0+i\sin y_0),
\quad x\in\mathbb R.
\]
\[
\text{The image is a ray w/ angle }y_0.
\]

\entry{Images of log}
\[
\log z \text{ transforms circles into vertical lines}
\quad\text{and}\quad
\text{rays into horizontal lines.}
\]

\topic{Complex Differentiation}

\entry{Analytic}
\[
f:U\to\mathbb C \text{ is analytic on an open }U\subseteq\mathbb C
\]
\[
\text{if } f \text{ is diff'ble } \forall u\in U
\text{ and } f' \text{ is continuous } \forall u\in U.
\]

\entry{Goursat's Thm}
\[
\text{If } f \text{ is diff'ble at every point of }D,\text{ then }f
\text{ is continuous on }D.
\]

\entry{Cauchy-Riemann Equations}
\[
f(z)=u+iv \text{ on }D\subseteq\mathbb C.
\]
\[
f \text{ is diff'ble at } z
\Longleftrightarrow
\text{C-R equations hold at }z:
\qquad u_x=v_y,\quad u_y=-v_x.
\]
\[
\Longleftrightarrow \text{the partials }u,v\text{ exist and are continuous on }\mathbb C.
\]

\entry{Constancy Thm}
\[
\text{Suppose }f\text{ is analytic on domain }D\subseteq\mathbb C.
\]
\[
(f'(z)=0\ \forall z\in D)\quad\text{or}\quad(|f(z)|\in\mathbb R\ \forall z\in D)
\quad\text{or}\quad(f,\overline f\text{ analytic on }D)
\]
\[
\Longrightarrow f \text{ is constant on }D.
\]

\newpage

\topic{Harmonic Functions}

\entry{Harmonic function}
\[
\text{Let }u:D\subseteq\mathbb R^2\to\mathbb R\text{ s.t. }u\in C^2.
\]
\[
\text{Say }u\text{ is harmonic on }D\text{ if }u_{xx}+u_{yy}=0\text{ on }D.
\]

\entry{Harmonic and Analytic Thm}
\[
\text{If }u+iv\text{ is analytic on }D,\text{ then }u\text{ and }v
\text{ are harmonic on }D.
\]

\entry{Harmonic Conjugate}
\[
\text{Let }u:D\subseteq\mathbb R^2\to\mathbb R\text{ be harmonic.}
\]
\[
\text{The harmonic conjugate }v:D\to\mathbb R
\text{ is a harmonic function for which }f=u+iv
\text{ is analytic on }D.
\]

\entry{Maybe tested: Existence Thm}
\[
\text{Let }u\text{ be harmonic on }D.\text{ If }D\text{ is simply connected,}
\]
\[
\text{then }\exists\text{ a harmonic conjugate }v\text{ on }D.
\]

\entry{Polygonally-connected}
\[
E\subseteq\mathbb C\text{ is polygonally-connected if any two points in }E
\]
\[
\text{can be connected by a sequence of line segments lying inside the set.}
\]

\entry{Simply connected}
\[
\text{A polygonally-connected region }D\text{ is simply-connected if}
\]
\[
\text{the region enclosed by any simple closed curve is contained in }D.
\]

\entry{Green's Thm}
\[
\text{If }D\subseteq\mathbb R^2\text{ is a bounded region with positively oriented boundary }\partial D,
\]
\[
\oint_{\partial D} P\,dx+Q\,dy
=
\iint_D (Q_x-P_y)\,dA.
\]

\newpage

\topic{Complex Integration}

\entry{Line Integrals}
\[
\text{Scalar line integral: }\quad
\int_\gamma f(x,y)\,ds
=
\int_a^b f(\gamma(t))\|\gamma'(t)\|\,dt.
\]
\[
\text{Vector line integral: }\quad
\int_\gamma \vec F\cdot d\vec r
=
\int_a^b F(\gamma(t))\cdot \gamma'(t)\,dt.
\]
\[
\text{Complex line integral: }\quad
\int_\gamma f(z)\,dz
=
\int_a^b f(\gamma(t))\gamma'(t)\,dt.
\]

\entry{ML estimate}
\[
\text{Let }\gamma\text{ be piecewise smooth, }f\text{ continuous on }\gamma.
\]
\[
\left|\int_\gamma f(z)\,dz\right|
\le
\int_\gamma |f(z)|\,|dz|.
\]
\[
\text{If }L=\ell(\gamma),\ |f(z)|\le M\text{ on }\gamma,\text{ then }
\left|\int_\gamma f(z)\,dz\right|\le ML.
\]

\entry{FTC}
\[
\text{I: Let }f\text{ be continuous on }D\subseteq\mathbb C,\quad F'(z)=f(z)\text{ on }D.
\]
\[
\text{Then for }\gamma\text{ w/ }\gamma(a)=A,\ \gamma(b)=B,\quad
\int_\gamma f(z)\,dz=\int_A^B f(z)\,dz=F(B)-F(A).
\]
\[
\text{II: Let }D\subseteq\mathbb C\text{ be simply connected, }f\text{ analytic on }D.
\]
\[
\text{Then }f\text{ has a primitive }F\text{ on }D,\qquad
F(z)=\int_{z_0}^z f(w)\,dw.
\]

\entry{Cauchy's Thm}
\[
\text{Let }D\text{ be a bounded domain with }\partial D
\text{ a finite number of disjoint contours.}
\]
\[
\text{If }f\text{ is analytic on }\partial D\cup D,\text{ then }
\int_{\partial D} f(z)\,dz=0.
\]

\entry{Morera's Thm}
\[
\text{Let }f\text{ be continuous on a domain }D.
\]
\[
\text{If }\int_c f(z)\,dz=0\text{ for every closed rectangle, contour, triangle in }D,
\]
\[
\text{then }f\text{ is analytic on }D.
\]

\entry{Cauchy Integral Formula}
\[
\text{Let }D\text{ be bounded w/ }\partial D\text{ a contour; if }f
\text{ is analytic on }\partial D\cup D,
\]
\[
\text{then } f(z_0)=\frac{1}{2\pi i}\int_{\partial D}\frac{f(z)}{z-z_0}\,dz,
\qquad z_0\in D.
\]
\[
\text{Generalized formula: }\quad
f^{(m)}(z_0)=\frac{m!}{2\pi i}\int_{\partial D}
\frac{f(z)}{(z-z_0)^{m+1}}\,dz,\qquad z_0\in D,\ m\ge0.
\]

\entry{Cauchy's Estimates}
\[
\text{Suppose }f\text{ is analytic for }|z-z_0|\le \rho.
\]
\[
\text{If }\exists M>0\text{ s.t. }|f(z)|\le M\ \forall z
\text{ w/ }|z-z_0|\le\rho,
\]
\[
\text{then }\left|f^{(m)}(z_0)\right|\le \frac{m!}{\rho^m}M.
\]

\entry{Liouville's Thm}
\[
\text{If }f(z)\text{ is entire and bounded, then }f(z)\text{ is constant.}
\]

\newpage

\topic{Complex Series}

\entry{Series convergence}
\[
\sum_{n=0}^{\infty}a_n\text{ conv}
\Longleftrightarrow
\lim_{N\to\infty}s_N=\lim_{N\to\infty}\sum_{n=0}^{N}a_n
\text{ converges.}
\]

\entry{Uniform convergence and analyticity}
\[
\text{If }f_n\to f\text{ uniformly on }D
\text{ and each }f_n\text{ is analytic on }D,
\]
\[
\text{then }f\text{ is analytic on }D.
\]

\entry{Taylor's Thm}
\[
\text{Suppose }f(z)\text{ is analytic for }|z-z_0|<R.
\]
\[
f(z)=\sum_{n=0}^{\infty}a_n(z-z_0)^n,\qquad
|z-z_0|<R,\qquad
a_n=\frac{f^{(n)}(z_0)}{n!}.
\]
\[
\text{Power series for }R<\infty:
\quad
\lim_{n\to\infty}\left|\frac{a_n}{a_{n+1}}\right|=R
\quad\text{or}\quad
\lim_{n\to\infty}\sqrt[n]{|a_n|}=\frac1R.
\]

\entry{Uniqueness principle}
\[
\text{If }f,g\text{ analytic on }D\text{ and }f(z)=g(z)
\text{ for all }z\in E,\text{ where }E\subseteq D
\]
\[
\text{contains a nonisolated point, then }f(z)=g(z)\quad \forall z\in D.
\]

\entry{Weierstrass M-test}
\[
\text{Suppose }M_n\ge0\text{ and }\sum M_n\text{ converges.}
\]
\[
\text{If }|g_n(z)|\le M_n\ \forall z\in E,\ \forall n\in\mathbb N,
\text{ then }\sum_{n=1}^{\infty}g_n(z)
\text{ conv unif and abs on }E.
\]

\entry{Cauchy product formula}
\[
\left(\sum a_nz^n\right)\left(\sum b_nz^n\right)
=
\sum c_nz^n,\qquad
c_n=\sum_{j=0}^n a_jb_{n-j}.
\]

\topic{Zeros and Singularities}

\entry{Zeros of function}
\[
\text{Suppose }f\text{ is analytic and }f\ne0\text{ on }D.
\]
\[
a\text{ zero of }f,\quad f(z)=(z-a)^nh(z)
\text{ for }h(a)\ne0.
\]
\[
a\text{ zero has order }n\text{ if }f(a)=\cdots=f^{(n-1)}(a)=0
\text{ and }f^{(n)}(a)\ne0.
\]

\entry{Laurent series}
\[
\text{If }f\text{ is analytic on annulus }\rho<|z-z_0|<R,
\]
\[
f(z)=\sum_{n=-\infty}^{\infty}a_n(z-z_0)^n.
\]
\[
\text{Analytic part: }\sum_{n=0}^{\infty}a_n(z-z_0)^n,\qquad
\text{principal part: }\sum_{n=-\infty}^{-1}a_n(z-z_0)^n.
\]

\entry{Isolated Singularities}
\[
z_0\text{ is an isolated singularity of }f
\text{ if }f\text{ is analytic on }D(z_0,r)\setminus\{z_0\}.
\]
\[
\text{Removable: }\lim_{z\to z_0}f(z)\text{ exists.}
\]
\[
\text{Pole of order }N:\quad f(z)=\frac{\phi(z)}{(z-z_0)^N},
\quad \phi\text{ analytic and }\phi(z_0)\ne0.
\]
\[
\text{Essential: not removable and not pole; Laurent principal part has infinitely many terms.}
\]

\entry{Quotient-pole lemma}
\[
\text{Suppose }f,g\text{ analytic in some disk centered at }z_0.
\]
\[
\text{If }f\text{ has a zero of order }m\text{ at }z_0
\text{ and }g\text{ has a zero of order }n\text{ at }z_0,
\]
\[
\frac{f(z)}{g(z)}\text{ has a removable singularity if }m\ge n,
\quad
\text{and a pole of order }n-m\text{ if }m<n.
\]

\newpage

\topic{Residue}

\entry{Residue}
\[
\text{Suppose }f\text{ has Laurent series expansion }
f(z)=\sum_{n=-\infty}^{\infty}a_n(z-z_0)^n
\text{ on }0<|z-z_0|<R.
\]
\[
\operatorname{Res}(f,z_0)=a_{-1}
=
\frac{1}{2\pi i}\int_\gamma f(z)\,dz.
\]

\entry{Residue Thm}
\[
\text{Let }D\text{ be a domain w/ piecewise smooth boundary }\partial D.
\]
\[
\text{Suppose }f\text{ is analytic on }D\cup\partial D
\text{ except at finite isolated singularities }z_1,\ldots,z_n,
\]
\[
\text{then }\int_{\partial D}f(z)\,dz
=2\pi i\sum_{j=1}^{n}\operatorname{Res}(f,z_j).
\]

\entry{Rules of Residue}
\[
\text{Rule 0: If }f\text{ has a removable singularity at }z_0,\text{ then }\operatorname{Res}(f,z_0)=0.
\]
\[
\text{Rule 1: If }f\text{ has a pole of order }1\text{ at }z_0,\text{ then }
\operatorname{Res}(f,z_0)=\lim_{z\to z_0}(z-z_0)f(z).
\]
\[
\text{Rule 2: If }f\text{ has a pole of order }N\text{ at }z_0,\text{ then }
\operatorname{Res}(f,z_0)=\lim_{z\to z_0}
\frac{1}{(N-1)!}\frac{d^{N-1}}{dz^{N-1}}\bigl((z-z_0)^Nf(z)\bigr).
\]
\[
\text{Rule 3: If }f,g\text{ are both analytic near }z_0
\text{ and }g\text{ has a simple zero at }z_0,
\]
\[
\text{then }\operatorname{Res}\left(\frac{f}{g},z_0\right)=\frac{f(z_0)}{g'(z_0)}.
\]

\entry{Application of Residue}
\[
\text{For }\int_{-\infty}^{\infty}\frac{P(x)}{Q(x)}\,dx
\text{ where }P,Q\text{ are polynomial and }\deg Q\ge \deg P+2,
\]
\[
\text{then }\int_{\partial D_R}\frac{P(z)}{Q(z)}\,dz
=
\int_{-R}^{R}\frac{P(x)}{Q(x)}\,dx
\;+\;\int_{\Gamma_R}\frac{P(z)}{Q(z)}\,dz.
\]
\[
\text{Evaluate by Residue Thm; arc integral goes to }0\text{ by ML estimate.}
\]

\end{document}
